% Options for packages loaded elsewhere
\PassOptionsToPackage{unicode}{hyperref}
\PassOptionsToPackage{hyphens}{url}
\PassOptionsToPackage{dvipsnames,svgnames,x11names}{xcolor}
%
\documentclass[
  11pt,
]{article}
\usepackage{amsmath,amssymb}
\usepackage{iftex}
\ifPDFTeX
  \usepackage[T1]{fontenc}
  \usepackage[utf8]{inputenc}
  \usepackage{textcomp} % provide euro and other symbols
\else % if luatex or xetex
  \usepackage{unicode-math} % this also loads fontspec
  \defaultfontfeatures{Scale=MatchLowercase}
  \defaultfontfeatures[\rmfamily]{Ligatures=TeX,Scale=1}
\fi
\usepackage{lmodern}
\ifPDFTeX\else
  % xetex/luatex font selection
\fi
% Use upquote if available, for straight quotes in verbatim environments
\IfFileExists{upquote.sty}{\usepackage{upquote}}{}
\IfFileExists{microtype.sty}{% use microtype if available
  \usepackage[]{microtype}
  \UseMicrotypeSet[protrusion]{basicmath} % disable protrusion for tt fonts
}{}
\makeatletter
\@ifundefined{KOMAClassName}{% if non-KOMA class
  \IfFileExists{parskip.sty}{%
    \usepackage{parskip}
  }{% else
    \setlength{\parindent}{0pt}
    \setlength{\parskip}{6pt plus 2pt minus 1pt}}
}{% if KOMA class
  \KOMAoptions{parskip=half}}
\makeatother
\usepackage{xcolor}
\usepackage[margin=25mm]{geometry}
\usepackage{color}
\usepackage{fancyvrb}
\newcommand{\VerbBar}{|}
\newcommand{\VERB}{\Verb[commandchars=\\\{\}]}
\DefineVerbatimEnvironment{Highlighting}{Verbatim}{commandchars=\\\{\}}
% Add ',fontsize=\small' for more characters per line
\newenvironment{Shaded}{}{}
\newcommand{\AlertTok}[1]{\textcolor[rgb]{1.00,0.00,0.00}{\textbf{#1}}}
\newcommand{\AnnotationTok}[1]{\textcolor[rgb]{0.38,0.63,0.69}{\textbf{\textit{#1}}}}
\newcommand{\AttributeTok}[1]{\textcolor[rgb]{0.49,0.56,0.16}{#1}}
\newcommand{\BaseNTok}[1]{\textcolor[rgb]{0.25,0.63,0.44}{#1}}
\newcommand{\BuiltInTok}[1]{\textcolor[rgb]{0.00,0.50,0.00}{#1}}
\newcommand{\CharTok}[1]{\textcolor[rgb]{0.25,0.44,0.63}{#1}}
\newcommand{\CommentTok}[1]{\textcolor[rgb]{0.38,0.63,0.69}{\textit{#1}}}
\newcommand{\CommentVarTok}[1]{\textcolor[rgb]{0.38,0.63,0.69}{\textbf{\textit{#1}}}}
\newcommand{\ConstantTok}[1]{\textcolor[rgb]{0.53,0.00,0.00}{#1}}
\newcommand{\ControlFlowTok}[1]{\textcolor[rgb]{0.00,0.44,0.13}{\textbf{#1}}}
\newcommand{\DataTypeTok}[1]{\textcolor[rgb]{0.56,0.13,0.00}{#1}}
\newcommand{\DecValTok}[1]{\textcolor[rgb]{0.25,0.63,0.44}{#1}}
\newcommand{\DocumentationTok}[1]{\textcolor[rgb]{0.73,0.13,0.13}{\textit{#1}}}
\newcommand{\ErrorTok}[1]{\textcolor[rgb]{1.00,0.00,0.00}{\textbf{#1}}}
\newcommand{\ExtensionTok}[1]{#1}
\newcommand{\FloatTok}[1]{\textcolor[rgb]{0.25,0.63,0.44}{#1}}
\newcommand{\FunctionTok}[1]{\textcolor[rgb]{0.02,0.16,0.49}{#1}}
\newcommand{\ImportTok}[1]{\textcolor[rgb]{0.00,0.50,0.00}{\textbf{#1}}}
\newcommand{\InformationTok}[1]{\textcolor[rgb]{0.38,0.63,0.69}{\textbf{\textit{#1}}}}
\newcommand{\KeywordTok}[1]{\textcolor[rgb]{0.00,0.44,0.13}{\textbf{#1}}}
\newcommand{\NormalTok}[1]{#1}
\newcommand{\OperatorTok}[1]{\textcolor[rgb]{0.40,0.40,0.40}{#1}}
\newcommand{\OtherTok}[1]{\textcolor[rgb]{0.00,0.44,0.13}{#1}}
\newcommand{\PreprocessorTok}[1]{\textcolor[rgb]{0.74,0.48,0.00}{#1}}
\newcommand{\RegionMarkerTok}[1]{#1}
\newcommand{\SpecialCharTok}[1]{\textcolor[rgb]{0.25,0.44,0.63}{#1}}
\newcommand{\SpecialStringTok}[1]{\textcolor[rgb]{0.73,0.40,0.53}{#1}}
\newcommand{\StringTok}[1]{\textcolor[rgb]{0.25,0.44,0.63}{#1}}
\newcommand{\VariableTok}[1]{\textcolor[rgb]{0.10,0.09,0.49}{#1}}
\newcommand{\VerbatimStringTok}[1]{\textcolor[rgb]{0.25,0.44,0.63}{#1}}
\newcommand{\WarningTok}[1]{\textcolor[rgb]{0.38,0.63,0.69}{\textbf{\textit{#1}}}}
\usepackage{longtable,booktabs,array}
\usepackage{calc} % for calculating minipage widths
% Correct order of tables after \paragraph or \subparagraph
\usepackage{etoolbox}
\makeatletter
\patchcmd\longtable{\par}{\if@noskipsec\mbox{}\fi\par}{}{}
\makeatother
% Allow footnotes in longtable head/foot
\IfFileExists{footnotehyper.sty}{\usepackage{footnotehyper}}{\usepackage{footnote}}
\makesavenoteenv{longtable}
\setlength{\emergencystretch}{3em} % prevent overfull lines
\providecommand{\tightlist}{%
  \setlength{\itemsep}{0pt}\setlength{\parskip}{0pt}}
\setcounter{secnumdepth}{-\maxdimen} % remove section numbering
\setlength{\emergencystretch}{3em}
\ifLuaTeX
  \usepackage{selnolig}  % disable illegal ligatures
\fi
\usepackage{bookmark}
\IfFileExists{xurl.sty}{\usepackage{xurl}}{} % add URL line breaks if available
\urlstyle{same}
\hypersetup{
  pdftitle={Infinite log-concavity of trinomial powers},
  colorlinks=true,
  linkcolor={blue},
  filecolor={Maroon},
  citecolor={Blue},
  urlcolor={blue},
  pdfcreator={LaTeX via pandoc}}

\title{Infinite log-concavity of trinomial powers}
\usepackage{etoolbox}
\makeatletter
\providecommand{\subtitle}[1]{% add subtitle to \maketitle
  \apptocmd{\@title}{\par {\large #1 \par}}{}{}
}
\makeatother
\subtitle{A proposed computer-assisted proof of the all-exponent
trinomial assertion}
\author{}
\date{31 October 2026}

\begin{document}
\maketitle

{
\hypersetup{linkcolor=}
\setcounter{tocdepth}{2}
\tableofcontents
}
\clearpage
\section{Abstract}\label{abstract}

Let \(a_{n,k}=[x^k](1+x+x^2)^n\), extended by zero outside its support,
and define \((\mathcal La)_k=a_k^2-a_{k-1}a_{k+1}\). We give a
computer-assisted proof that \(\mathcal L^r a_n\) is entrywise
nonnegative for every integer \(n\ge10\) and every integer \(r\ge0\).
The proof joins an exact invariant-cone verification for
\(10\le n\le4003\) with a uniform analytic and algebraic argument for
\(n\ge4004\). In the latter range, a sixteen-variable integral and
explicit saddle estimates supply bulk signs for the fourth iterate;
exact moving-node certificates supply the remaining edge signs.
Palindromicity and a degree count then prove negative real-rootedness of
the fourth iterate. Separate Toeplitz-minor arguments establish the
earlier layers, and a real-rootedness preservation theorem gives every
later layer.

The numerical components use exact integers and rational enclosures. The
claimed quantifiers are not inferred from a finite sample of exponents.
This draft has undergone independent internal checks, including
separately implemented arithmetic checks; it has not undergone external
peer review or proof-assistant formalization. No assertion of exhaustive
literature priority is made.

\section{The target and the proof
architecture}\label{the-target-and-the-proof-architecture}

\textbf{Theorem.} For every integer \(n\ge10\), every integer \(r\ge0\),
and every integer \(k\), \[
 (\mathcal L^r a_n)_k\ge0,
 \qquad a_{n,k}=[x^k](1+x+x^2)^n.
\] All sequences and iterates use zero extension beyond their natural
finite support.

This is the specific trinomial assertion discussed in Example 5.3 of
Medina and Straub, \emph{On multiple and infinite log-concavity},
arXiv:1405.1765v1, printed page 11; the published article appears in
\emph{Annals of Combinatorics} \textbf{20} (2016), 125--138. The
original work verified a finite range \(10\le n\le500\). The present
theorem does not assert the more general polynomial conjecture
considered elsewhere in that paper.

The proof has four logically separate components:

\begin{enumerate}
\def\labelenumi{\arabic{enumi}.}
\tightlist
\item
  \textbf{Finite exponents.} Every \(10\le n\le4003\) reaches a
  sufficient invariant cone under exact directed ratio arithmetic. The
  cone proves all subsequent iterations at once.
\item
  \textbf{Earlier layers in the infinite range.} Finite certificates at
  exponents eight and nine imply the full \(\mathrm{PF}_4\) property
  through symmetric-function identities. Convolution closure and
  determinant identities then prove nonnegativity through the third
  iterate for every \(n\ge56\).
\item
  \textbf{Fourth-iterate roots.} Uniform bulk estimates and exact
  certificates on the continuous parameter interval \(0\le u\le1/4004\)
  provide enough alternating real signs to locate all \(2n\) zeros of
  the fourth iterate on the negative real axis for every integer
  \(n\ge4004\).
\item
  \textbf{Later layers.} Brändén's preservation theorem applies to the
  fourth iterate and thereafter. It is not used backwards to infer
  positivity of the earlier layers.
\end{enumerate}

The finite and infinite exponent ranges are complementary and leave no
gap. The roles of the continuous parameter family and the finite
polynomial at reciprocal integers are kept distinct throughout.

\section{Verification conventions}\label{verification-conventions}

A rational enclosure is a proved inclusion, not a floating-point
approximation. Every strict sign assertion subtracts all applicable
rounding and analytic remainder bounds before it is accepted. The exact
arithmetic backend is credited and pinned in the reproducibility
package. A fast replay of saved enclosures and a full regeneration from
the input formulas have different trust boundaries; the package
documents both rather than conflating them.

The mathematical statements below are accompanied by portable
verification commands and explicit certificate coverage. Raw development
logs and internal review narratives are not part of the mathematical
exposition.

\section{Finite certificates and the first three
iterates}\label{finite-certificates-and-the-first-three-iterates}

Let \(p(x)=1+x+x^2\) and \(a_k=[x^k]p(x)^n\), with zero extension. All
arithmetic certificates below use integers and directed rational
enclosures.

\subsection{\texorpdfstring{1. Exact ratio verification for
\(10\le n\le4003\)}{1. Exact ratio verification for 10\textbackslash le n\textbackslash le4003}}\label{exact-ratio-verification-for-10le-nle4003}

The coefficient identity \(p(x)(p(x)^n)'=np'(x)p(x)^n\) gives \[
(k+1)a_{k+1}=(n-k)a_k+(2n-k+1)a_{k-1}.
\] The coefficients are positive on \(0\le k\le2n\) and symmetric. On
the first half define \(r_k=a_k/a_{k-1}\). Then \(r_1=n\) and \[
r_{k+1}=\frac{n-k+(2n-k+1)/r_k}{k+1},\qquad 1\le k<n.
\] This expression decreases with \(r_k\). If \(S=2^{256}\) and
\(R_k^-\le Sr_k\le R_k^+\), directed integer divisions give \[
\begin{aligned}
R_{k+1}^-&=\left\lfloor\frac{(n-k)SR_k^++(2n-k+1)S^2}{(k+1)R_k^+}\right\rfloor,\\
R_{k+1}^+&=\left\lceil\frac{(n-k)SR_k^-+(2n-k+1)S^2}{(k+1)R_k^-}\right\rceil.
\end{aligned}
\] All denominators are positive, as checked by the certificate program.

For \(q_k=a_{k-1}a_{k+1}/a_k^2\), the initial half-row enclosures are \[
\begin{aligned}
Q_k^-&=\left\lfloor\frac{SR_{k+1}^-}{R_k^+}\right\rfloor,&
Q_k^+&=\left\lceil\frac{SR_{k+1}^+}{R_k^-}\right\rceil,
&&1\le k<n,\\
Q_n^-&=\left\lfloor\frac{S^3}{(R_n^+)^2}\right\rfloor,&
Q_n^+&=\left\lceil\frac{S^3}{(R_n^-)^2}\right\rceil,
\end{aligned}
\] and \(Q_0^-=Q_0^+=0\). Symmetry gives the remaining indices. The
powers of \(S\) in the central formula are essential.

Whenever \(0\le q_k<1\), the next iterate is positive on its support and
its ratios obey \[
q'_k=\frac{q_k^2(1-q_{k-1})(1-q_{k+1})}{(1-q_k)^2}.
\] The map increases in its own argument and decreases in either
neighboring argument. Given integer enclosures
\(0\le Q^-\le Sq\le Q^+<S\), lower and upper bounds for \(Sq'_k\) are
therefore \[
\left\lfloor\frac{(Q_k^-)^2(S-Q_{k-1}^+)(S-Q_{k+1}^+)}{S(S-Q_k^-)^2}\right\rfloor,
\qquad
\left\lceil\frac{(Q_k^+)^2(S-Q_{k-1}^-)(S-Q_{k+1}^-)}{S(S-Q_k^+)^2}\right\rceil.
\] At the center of the symmetric half-array the right neighbor equals
the left neighbor. No unbounded integer is replaced by floating-point
arithmetic.

The exact program checks every integer \(n=10,\ldots,4003\). At each
step it verifies the strict upper bound \(Q_k^+<S\), and it terminates
only when \(3Q_k^+\le S\) for every \(k\). All 3994 cases terminate by
depth 12. A separate verifier using exact coefficient rows instead of
the successive-ratio initializer gives the same successful coverage. The
terminal condition is sufficient because \(q\le1/3\) implies \[
q'\le\frac{(1/3)^2}{(1-1/3)^2}=\frac14<\frac13.
\] Induction proves positivity of all later iterates, so a finite exact
certificate per exponent establishes infinite log-concavity rather than
merely a bounded number of successful iterations.

\subsection{\texorpdfstring{2. Finite certificates for the full
\(\mathrm{PF}_4\)
property}{2. Finite certificates for the full \textbackslash mathrm\{PF\}\_4 property}}\label{finite-certificates-for-the-full-mathrmpf_4-property}

For a zero-extended finite sequence \(a\), let \(T(a)_{i,j}=a_{j-i}\),
\(i,j\ge0\). The property \(\mathrm{PF}_4\) means every minor of order
at most four is nonnegative.

Suppose \[
A(x)=\sum_{k=0}^d a_kx^k=\prod_{h=1}^d(1+z_hx),
\] so \(a_k=e_k(z)\). It suffices to test the finite determinants \[
\det(a_{u_i-i+j})_{1\le i,j\le4}
\] for every \(d\ge u_1\ge u_2\ge u_3\ge u_4\ge0\). Here the \(z_h\) may
be complex; the tests concern real evaluated determinants.

Indeed dual Jacobi--Trudi identifies these determinants with ordinary
Schur functions \(s_\lambda(z)\) for \(\lambda'=u\), namely all shapes
of width at most four and height at most \(d\); zero padding includes
smaller widths. Consider an arbitrary order-\(\ell\le4\) Toeplitz minor
with increasing rows \(r_i\) and columns \(c_i\). If \(c_i<r_i\) for
some \(i\), the first \(i\) columns can meet at most \(i-1\) eligible
rows, so every determinant term vanishes. Otherwise put \[
\lambda'_i=c_{\ell+1-i}+i,\qquad \mu'_i=r_{\ell+1-i}+i.
\] These are partitions with \(\lambda'\ge\mu'\). Transposition and
simultaneous reversal transform the minor into \[
\det(e_{\lambda'_i-\mu'_j-i+j})=s_{\lambda/\mu}(z).
\] The Littlewood--Richardson expansion \[
s_{\lambda/\mu}=\sum_\nu c_{\mu\nu}^{\lambda}s_\nu
\] has nonnegative integer coefficients. For nonzero terms
\(\nu\subseteq\lambda\), hence \(\operatorname{width}(\nu)\le\ell\le4\).
If \(\operatorname{height}(\nu)>d\), its specialization to \(d\)
variables vanishes. Every remaining ordinary Schur value is among the
finite tests. Thus every arbitrary Toeplitz minor is nonnegative.

For \(p(x)^n\) use \(d=2n\) and \(n\) copies each of \(\exp(i\pi/3)\)
and \(\exp(-i\pi/3)\). The exact tests at \(n=8\) and \(n=9\) enumerate
respectively \(\binom{20}{4}=4845\) and \(\binom{22}{4}=7315\)
determinants. Every determinant is positive; the minimum is 1. These are
finite certificates for the complete infinite-matrix \(\mathrm{PF}_4\)
condition, not merely a sample of consecutive-row minors.

\subsection{3. Convolution closure and nonnegative early
iterates}\label{convolution-closure-and-nonnegative-early-iterates}

We have \(T(a*b)=T(a)T(b)\). Each fixed minor of the product has a
finite Cauchy--Binet sum: contributing intermediate indices lie between
the selected row and column indices. If both factors are
\(\mathrm{PF}_4\), every summand is nonnegative. Thus convolution
preserves \(\mathrm{PF}_4\).

For \(n\ge56\) write \(n=8q+r\), \(0\le r\le7\). Since \(q\ge7\ge r\),
\[
n=8(q-r)+9r.
\] Hence \(p^n\) is a convolution of the two certified \(\mathrm{PF}_4\)
seed rows, and is \(\mathrm{PF}_4\).

Let \(D_j(k)=\det(a_{k+s-t})_{0\le s,t<j}\), \(D_0=1\). Desnanot--Jacobi
gives \[
D_j(k)^2-D_j(k-1)D_j(k+1)=D_{j-1}(k)D_{j+1}(k).
\] It follows algebraically that \[
\begin{aligned}
(\mathcal La)_k&=D_2(k),\\
(\mathcal L^2a)_k&=a_kD_3(k),\\
(\mathcal L^3a)_k&=D_2(k)\bigl[a_k^2D_4(k)+D_3(k-1)D_3(k+1)\bigr].
\end{aligned}
\] These are polynomial identities with no division by a minor, so
endpoints and possible zero minors cause no exception. The property
\(\mathrm{PF}_4\) therefore implies nonnegativity of layers zero through
three. The selected infinite range \(n\ge4004\) lies entirely within
this proven range \(n\ge56\).

\subsection{References}\label{references}

The symmetric-function identities are dual Jacobi--Trudi and
Littlewood--Richardson; see I. G. Macdonald, \emph{Symmetric Functions
and Hall Polynomials}, second edition, Chapter I, Sections 3, 5 and 9.
The first three iteration identities and fourth-order limitations are
discussed by P. R. W. McNamara and B. E. Sagan, \emph{Infinite
log-concavity: developments and conjectures},
\url{https://arxiv.org/abs/0808.1065}. These general algebraic tools are
prior results, not new claims of this proof.

\clearpage

\section{\texorpdfstring{Analytic part of a computer-assisted proof for
powers of
\(1+x+x^2\)}{Analytic part of a computer-assisted proof for powers of 1+x+x\^{}2}}\label{analytic-part-of-a-computer-assisted-proof-for-powers-of-1xx2}

\subsection{1. Statement and structure of the
proof}\label{statement-and-structure-of-the-proof}

For a finitely supported sequence \(a=(a_k)_{k\in\mathbb Z}\), extended
by zero outside its support, write \[
 (\mathcal La)_k=a_k^2-a_{k-1}a_{k+1}.
\] The sequence is \textbf{infinitely log-concave} if every iterate
\(\mathcal L^j a\) is nonnegative.

The result addressed by this draft is the following integer-exponent
statement:

\begin{quote}
\textbf{Theorem.} For every integer \(n\ge10\), the coefficient sequence
of \((1+x+x^2)^n\) is infinitely log-concave.
\end{quote}

This is the trinomial case discussed in Medina--Straub, Example 5.3 and
Table 1. It is not the general assertion for every polynomial with
positive coefficients in their Conjecture 5.2. See
\href{https://arminstraub.com/downloads/pub/infinitelogconcavity.pdf}{Medina--Straub,
\emph{On multiple and infinite log-concavity}}.

The proof is computer-assisted in two explicit ways: certain low-degree
rational polynomial inequalities are verified by positive Bernstein
coefficients, and a finite collection of parameter-wide sign
inequalities is verified by exact interval Taylor arithmetic.
Floating-point calculations may select rational evaluation points, but
do not certify their signs.

Put \[
 p(z)=1+z+z^2,\qquad a_{n,k}=[z^k]p(z)^n,
 \qquad f_{n,j}(x)=\sum_k(\mathcal L^j a_n)_k x^k.
\] The analytic argument proves that \(f_{n,4}\) has only negative real
zeros for every integer \(n\ge4004\). Its degree is \(2n\). We do not
require, or assert here, this stronger real-rootedness statement
throughout the smaller finite range.

The proof has three parts.

\begin{enumerate}
\def\labelenumi{\arabic{enumi}.}
\tightlist
\item
  A uniform two-saddle estimate gives alternating signs of \(f_{n,4}\)
  at explicitly defined bulk phase points.
\item
  Exact parameter-wide certificates supply a finite chain of edge signs
  and its connection to the bulk.
\item
  The intermediate value theorem and reciprocity exhaust the degree. The
  initial three iterates and the finite range \(10\le n\le4003\) are
  treated in the companion finite and algebraic chapter above.
\end{enumerate}

Once \(f_{n,4}\) has negative real zeros, Brändén's preservation theorem
applies repeatedly to the same zero-extended operator \(\mathcal L\).
The theorem and its proof are in §3 of
\href{https://arxiv.org/pdf/0909.1927}{Brändén, \emph{Iterated sequences
and the geometry of zeros}}. It is used forward from the fourth iterate,
not backward to infer the first three iterates.

Throughout, \(\operatorname{sinc}x=\sin(x)/x\), with its continuous
value at zero. Arguments of analytic functions are chosen continuously
from the positive real axis whenever this is specified.

\subsection{2. An exact sixteen-leaf
integral}\label{an-exact-sixteen-leaf-integral}

\subsubsection{2.1. The diagonal identity}\label{the-diagonal-identity}

Consider a complete balanced binary tree with \(r=2^j\) labelled leaves.
For an internal node \(\nu\), let \(L_\nu,R_\nu\) be its two child leaf
sets, and put \[
 Z_{\nu,L}=\prod_{i\in L_\nu}z_i,
 \qquad Z_{\nu,R}=\prod_{i\in R_\nu}z_i.
\] Define \[
 P_j(z)=\prod_{\nu}\left(1-\frac{Z_{\nu,L}}{Z_{\nu,R}}\right).
\] Then \[
 f_{n,j}(x)=\operatorname{CT}_{z_1,\ldots,z_{r-1}}
 P_j(z)\prod_{i=1}^r p(z_i)^n,
 \qquad z_r=\frac{x}{z_1\cdots z_{r-1}}.                 \tag{2.1}
\] To prove this, start with the coefficientwise product formula for a
diagonal. At one application of \(\mathcal L\), the product of two
equal-index coefficients is replaced by that product minus the product
with indices shifted by \(-1,+1\). These shifts insert exactly
\(Z_{\nu,L}/Z_{\nu,R}\). Induction down the binary tree proves (2.1),
including zero extension at the boundaries.

For \(x=-s^{16}\), take all sixteen radii equal to \(s\), and write
\(z_i=s e^{i\alpha_i}\). The angular constraint is \[
 \sum_{i=1}^{16}\alpha_i=\pi\pmod {2\pi}.
\] The measure is normalized Haar measure on this fifteen-dimensional
torus. In the coordinates \(\alpha_1,\ldots,\alpha_{15}\), it is
\((2\pi)^{-15}d\alpha_1\cdots d\alpha_{15}\).

\subsubsection{2.2. A nonnegative kernel}\label{a-nonnegative-kernel}

For every internal node set \[
 D_\nu(\alpha)=\sum_{i\in L_\nu}\alpha_i-
                    \sum_{i\in R_\nu}\alpha_i.
\] The exponential factor \(\prod_i p(s e^{i\alpha_i})^n\) and the
constrained measure are invariant under permutations of the leaves.
Average the kernel over the group generated by interchanging the two
children at each tree node. Beginning at the lowest level, an
interchange changes \(D_\nu\) to \(-D_\nu\), while the already averaged
descendant factors are even and merely permuted. It therefore replaces
\(1-e^{iD_\nu}\) by \(1-\cos D_\nu\). Repeating upward gives the exact
integral identity \[
 f_{n,4}(-s^{16})=
 \int_{\sum\alpha_i=\pi\ (2\pi)}
 K(\alpha)\prod_{i=1}^{16}p(s e^{i\alpha_i})^n\,d\mu(\alpha),
 \qquad K=\prod_\nu(1-\cos D_\nu)\ge0.                \tag{2.2}
\] The replacement is an equality after integration; it is not a
pointwise equality with the original Laurent kernel.

At a configuration of equal angles, the leading homogeneous term of
\(K\) is \[
 H(\eta)=\prod_\nu\frac{D_\nu(\eta)^2}{2}.           \tag{2.3}
\] There are fifteen internal nodes, so this vanishing has degree
thirty.

\subsubsection{2.3. Haar coordinates and measure
normalization}\label{haar-coordinates-and-measure-normalization}

On the tangent space \[
 \mathcal T=\{\eta\in\mathbb R^{16}:\sum_i\eta_i=0\},
\] the fifteen contrast vectors defining \(D_\nu\) are orthogonal. Their
squared norms are the node sizes \(m_\nu\). Sizes \(2,4,8,16\) occur
respectively \(8,4,2,1\) times. For orthonormal coordinates
\(x_\nu=D_\nu/\sqrt{m_\nu}\), \[
 H=\left(\prod_\nu\frac{m_\nu}{2}\right)\prod_\nu x_\nu^2
   =2048\prod_\nu x_\nu^2.                            \tag{2.4}
\] The original fifteen angular coordinates have metric matrix
\(I+\mathbf1\mathbf1^T\), of determinant sixteen. Thus their measure is
one quarter of Euclidean measure on \(\mathcal T\). This factor is
important in every saddle and complement estimate below.

\subsection{3. Saddle geometry and explicit scalar
inequalities}\label{saddle-geometry-and-explicit-scalar-inequalities}

Set \[
 \theta=\frac\pi{16},\quad \zeta=s e^{i\theta},\quad
 q(z)=1+4z+z^2,\quad
 m(z)=\frac{zp'(z)}{p(z)},\quad
 v(z)=\frac{zq(z)}{p(z)^2},\quad V=n v(s).
\] For \(m\ge2\), define \[
 \kappa_m(z)=\left(z\frac d{dz}\right)^{m-2}v(z).
\] Here \(v(s)>0\) for \(s>0\). We work with \(0<s\le1\); reciprocity
later supplies the other half of the zeros.

\subsubsection{3.1. Modulus estimates}\label{modulus-estimates}

Write \[
 h_s(\alpha)=\log\frac{|p(s e^{i\alpha})|}{p(s)}.
\] If \(c=\cos\alpha\), the exact identity is \[
 \frac{|p(s e^{i\alpha})|^2}{p(s)^2}=1-D(c),
 \quad D(c)=2a(1-c)+b(1-c^2),
\] where \[
 a=\frac{s(1+s^2)}{p(s)^2},\qquad
 b=\frac{4s^2}{p(s)^2},\qquad a+b=v(s),\qquad b\le\frac23v(s).
                                                               \tag{3.1}
\] In particular, if \(|\alpha|\ge\pi/3\), \[
 h_s(\alpha)\le-\frac5{12}v(s).                       \tag{3.2}
\] A small-angle bound at the equal-phase saddle is \[
 16h_s(\theta)>-\frac{40}{123}v(s).                  \tag{3.3}
\] For example, \(D(\cos\theta)\le v(s)\theta^2\), \(v(s)\le1\), and
\(-\log(1-x)\le x/(1-x)\) give \[
 -16h_s(\theta)\le
 \frac{\pi^2}{32(1-\pi^2/256)}v(s)<\frac{40}{123}v(s).
\]

On the central interval, \[
 \Re v(s e^{i\alpha})\ge\frac13v(s),
 \qquad |\alpha|\le\frac\pi3.                       \tag{3.4}
\] Consequently \(h_s''(\alpha)=-\Re v(s e^{i\alpha})\le-v(s)/3\): the
function is strictly concave there.

These facts locate the two modulus maxima. A configuration containing a
noncentral angle has total modulus loss at least \(5v(s)/12\), larger
than the equal-phase loss in (3.3). If all angles are central, the
possible winding sums are \(\pm\pi,\pm3\pi,\pm5\pi\). Jensen's
inequality shows that the largest modulus in each such sector occurs at
its equal-angle configuration. Since \(h_s\) is even and strictly
decreasing on \((0,\pi/3]\), the global maxima are exactly the two
configurations with all angles equal to \(\theta\) or \(-\theta\).

\subsubsection{3.2. Scalar certificate
lemma}\label{scalar-certificate-lemma}

The following bounds hold uniformly for \(0<s\le1\). Put \(v_0=v(s)\),
\(\beta=\Re v(\zeta)\). \[
 \frac{19}{20}v_0\le\beta\le\frac{17}{16}v_0,
 \qquad \left(\frac{|v(\zeta)|}{\beta}\right)^{45/2}<2,             \tag{3.5}
\] \[
 |\Im\kappa_3(\zeta)|\le\frac14v_0,
 \quad |\kappa_3(\zeta)|\le\frac65v_0,
 \quad |\kappa_4(\zeta)|\le\frac43v_0,
 \quad |\kappa_5(\zeta)|\le\frac94v_0.                            \tag{3.6}
\] On the wider sector \(\cos\alpha\ge3/4\), \[
 \Re\kappa_4(s e^{i\alpha})\le\frac32v_0,
 \quad |\kappa_4(s e^{i\alpha})|\le3v_0,
 \quad |\kappa_6(s e^{i\alpha})|\le32v_0.                        \tag{3.7}
\]

These are rational polynomial inequalities, rather than numerical bounds
on a grid. Here is an explicit way to verify them. If
\(\kappa_m(z)=N_m(z)/p(z)^m\), then \[
 N_2=zq(z),\qquad N_{m+1}=z\{pN_m'-m p'N_m\}.                   \tag{3.8}
\] For a polynomial \(N(z)=\sum a_i z^i\), \[
 |N(s e^{i\alpha})|^2
 =\sum_{i,j}a_i a_j s^{i+j}T_{|i-j|}(\cos\alpha),               \tag{3.9}
\] where \(T_j\) is a Chebyshev polynomial. Clearing positive
denominators in each squared modulus bound produces a polynomial in
\(s,c\). On the narrow sector, use \(c\in[49/50,1]\), containing
\(\cos\theta\), and split \(s\in[0,1]\) into four equal intervals. On
the wide sector use \(c\in[3/4,1]\). Nonnegative tensor-product
Bernstein controls verify the required inequalities. For a power-basis
polynomial \(\sum a_{ij}x^iy^j\) of bidegree \((d,e)\), the controls are
explicitly \[
 b_{kl}=\sum_{i\le k,\,j\le l}
 a_{ij}\frac{\binom{k}{i}}{\binom d i}
        \frac{\binom l j}{\binom e j}.                           \tag{3.10}
\] Thus the verification uses rational arithmetic only.

For completeness, some especially useful reductions are small. Let
\(t=s+s^{-1}\), \(c=\cos\alpha\), \[
 A=tc+1,\quad B^2=(t^2-4)(1-c^2),\quad D=A^2+B^2,
 \quad R=(A+3)A^2+(A-3)B^2.
\] Then \(\Re v=R/D^2\) and \(v(s)=(t+4)/(t+1)^2\). The polynomial \[
 R(t+1)^2-\tfrac13(t+4)D^2
\] has positive Bernstein controls after \(t=2+u\), \(c=(1+y)/2\): the
minima by powers \(u^0,\ldots,u^5\) are \[
 148,\quad347/3,\quad65/3,\quad1/2,\quad2/3,\quad1/6.
\] This proves (3.4). The conditioning inequality follows by verifying
\(16(\Im v)^2<(\Re v)^2\); hence
\((|v|/\Re v)^{45/2}<(17/16)^{45/4}<2\). The lower real-part bound in
(3.5) and the bounds involving real or imaginary parts in (3.6)--(3.7)
are verified after the same denominator clearing, using \(t\ge2\) where
convenient. The upper real-part bound also follows directly by
integrating \(|\kappa_4|\le3v_0\) twice from angle zero: \[
 \beta\le v_0(1+3\theta^2/2)<17v_0/16.
\] The accompanying exact certificate routines specify all polynomial
subdivisions; their role is confined to these finite rational
inequalities.

\subsection{4. Leading amplitude and first
correction}\label{leading-amplitude-and-first-correction}

Near the positive saddle write \(\alpha_i=\theta+\eta_i\),
\(\eta\in\mathcal T\). The exponent has quadratic part \[
 -\frac{n v(\zeta)}2\sum_i\eta_i^2.
\] Using (2.4), the complex Gaussian integral is a product of fifteen
one-dimensional second-moment integrals. Including the factor \(1/4\)
from the measure gives the one-saddle amplitude \[
 A_n(s)=\frac{2048}{4(2\pi)^{15/2}}
        \frac{p(\zeta)^{16n}}{(n v(\zeta))^{45/2}}.              \tag{4.1}
\] The square root and power are chosen continuously with
\(\Re v(\zeta)>0\). The negative saddle contributes the complex
conjugate.

Let \(S_j=\sum_i\eta_i^j\) and \(Q=\sum_\nu D_\nu^2\). The first
correction to the Gaussian integrand is \[
 -\frac Q{12}+\frac{n\kappa_4(\zeta)}{24}S_4
       -\frac{n^2\kappa_3(\zeta)^2}{72}S_3^2.                 \tag{4.2}
\] The relevant normalized unit-scale weighted Gaussian moments are \[
 \mathbb EQ=192,\qquad \mathbb ES_4=5055/16,
 \qquad \mathbb ES_3^2=2835/4.
\] They yield \[
 B_1(s)=-\frac{16}{v(\zeta)}
  +\frac{1685\kappa_4(\zeta)}{128v(\zeta)^2}
  -\frac{315\kappa_3(\zeta)^2}{32v(\zeta)^3}.                 \tag{4.3}
\] Equations (3.5)--(3.6) imply \[
 \left|\frac{B_1(s)}n\right|
 \le\frac{2173895}{41154V}<\frac{53}{V}.                     \tag{4.4}
\] No approximation to the zeros is inferred from this formal correction
alone. The next two sections supply actual remainder bounds.

\subsection{5. A uniform local
remainder}\label{a-uniform-local-remainder}

Let \[
 \mathcal B=\{\eta\in\mathcal T:|\eta_i|\le2/5\text{ for all }i\}.
\] We compare the exact integral over this box to the Gaussian plus
(4.2) over the \textbf{same} box. The resulting one-saddle error is \[
 \left|I_{\mathcal B}-I_{\mathcal B}^{(1)}\right|
 \le \frac{36821}{V^2}|A_n(s)|,\qquad V\ge100.                 \tag{5.1}
\] Here the factors \(p(\zeta)^{16n}\) and angular normalization are
included in both integrals.

\subsubsection{5.1. An even homotopy}\label{an-even-homotopy}

Put \[
 L(\eta)=\sum_i\{\log p(\zeta e^{i\eta_i})-\log p(\zeta)\},
 \quad E_\tau=nL(\tau\eta)/\tau^2,
\] \[
 F_\tau=\prod_\nu\operatorname{sinc}^2(\tau D_\nu/2),
 \qquad J(\tau,\eta)=H F_\tau e^{E_\tau}.
\] All singularities at \(\tau=0\) are removable. In particular,
\(E_0=-nv(\zeta)\|\eta\|^2/2\). Since \(J(-\tau,\eta)=J(\tau,-\eta)\),
integration over the symmetric box is even in \(\tau\). Taylor's
integral remainder therefore gives exactly \[
 \int_{\mathcal B}\{J(1)-J(0)-\tfrac12J''(0)\}
 =\int_0^1\frac{(1-\tau)^3}{6}
       \int_{\mathcal B}\partial_\tau^4J(\tau,\eta)\,d\eta\,d\tau.
                                                               \tag{5.2}
\] Parity is used only in this integrated identity. Odd terms in the
fourth derivative cannot be discarded at nonzero \(\tau\), because
\(e^{E_\tau}\) need not be even in \(\eta\).

\subsubsection{5.2. Bounds for the fourth
derivative}\label{bounds-for-the-fourth-derivative}

For \(m\ge1\), Taylor's integral formula gives \[
 \partial_\tau^m E_\tau
 =n\sum_i(i\eta_i)^{m+2}
       \int_0^1 u^m(1-u)
       \kappa_{m+2}(\zeta e^{iu\tau\eta_i})\,du.                \tag{5.3}
\] Set \(d=2/5\), \(A=|S_3|\), \(B=S_4\), \(C=|S_5|\), \(D_6=S_6\). The
scalar bounds imply \[
\begin{aligned}
 e_1&=V\{A/5+(\tau/9+9d\tau^2/160+8d^2\tau^3/45)B\},\\
 e_2&=V(1/9+9d\tau/80+8d^2\tau^2/15)B,\\
 e_3&=V(9C/80+16\tau D_6/15),\\
 e_4&=16VD_6/15,
\end{aligned}                                                   \tag{5.4}
\] with \(|\partial_\tau^mE_\tau|\le e_m\). In the first two lines only,
use \(|S_5|\le dS_4\), \(S_6\le d^2S_4\).

The identity \(\operatorname{sinc}^2(t/2)=\mathbb E e^{it(U+U')}\), for
independent uniforms on \([-1/2,1/2]\), gives \[
 |F|\le1,\quad |F'|\le\tau Q/6=f_1,\quad |F''|\le Q/6=f_2,
\] \[
 |F'''|\le\tau Q^2/12=f_3,\qquad |F''''|\le Q^2/12=f_4.         \tag{5.5}
\] For example, the exact fourth moment is \(Q^2/12-\sum D_\nu^4/60\),
so dropping the last term gives an upper bound. The odd derivative
estimates follow by integrating the next derivative from zero.

The product rule now gives \[
 |J^{(4)}|\le H e^{\Re E_\tau}\mathcal P,
\] \[
\begin{aligned}
\mathcal P={}&f_4+4f_3e_1+6f_2(e_2+e_1^2)
 +4f_1(e_3+3e_1e_2+e_1^3)\\
 &+e_4+4e_1e_3+3e_2^2+6e_1^2e_2+e_1^4.                       \end{aligned} \tag{5.6}
\] This polynomial has nonnegative coefficients in its nonnegative
arguments.

\subsubsection{5.3. Gaussian domination and a reproducible finite moment
sum}\label{gaussian-domination-and-a-reproducible-finite-moment-sum}

Real-part Taylor bounds give \[
 \Re E_\tau\le-\frac n2a\|\eta\|^2,
 \qquad a=\beta-\frac4{75}v_0\ge\frac{269}{300}v_0
                         >\left(\frac{15}{16}\right)^2v_0.
\] The ratio of this real Gaussian normalization to the absolute complex
Gaussian normalization is \[
 (|v(\zeta)|/a)^{45/2}<2(285/269)^{45/2}<8.                  \tag{5.7}
\] Under the normalized \(H\)-weighted Gaussian at unit scale, the
orthonormal Haar coordinates are independent with density
\(x^2e^{-x^2/2}/\sqrt{2\pi}\). Hence \[
 \mathbb ED_\nu^{2r}=(2r+1)!!m_\nu^r,
 \qquad \mathbb ED_\nu^{2r+1}=0.                              \tag{5.8}
\] The reconstruction \[
 \eta_i=\sum_{\nu\text{ ancestor of }i}
       \frac{\epsilon_{i,\nu}D_\nu}{m_\nu},\qquad
 \epsilon_{i,\nu}\in\{-1,1\},                                \tag{5.9}
\] makes every moment needed in (5.6) a finite rational sum. Absolute
odd powers of \(S_3\) or \(S_5\) are bounded by Cauchy--Schwarz, moving
one odd factor to each side so that both resulting powers are even. An
integer ceiling of the resulting square root is used. For example, \[
 \mathbb ES_3^4=38064735/16,\quad
 \mathbb ES_5^2=81661725/256,\quad
 \mathbb ES_4^4=2166032330330625/65536.
\]

For a monomial \(C_0\tau^pV^lA^aB^bC^cD_6^dQ^e\) in (5.6), let
\(q=3a+4b+5c+6d+2e\) denote its angular degree. Every monomial has
\(q-2l\ge4\). Its contribution to the coefficient of \(V^{-2}\), for
\(V\ge100\), is bounded by \[
 8\,\frac{p!}{(p+4)!}\,
 [\text{monomial coefficient}]\,[\text{unit moment bound}]
 \left(\frac{16}{15}\right)^q10^{2l-q+4}.                    \tag{5.10}
\] Here \(p!/(p+4)!\) is the exact integral of \(\tau^p(1-\tau)^3/6\);
replacing all these weights by \(1/24\) is unnecessarily wasteful.

Expanding the explicit polynomial (5.6), applying (5.8)--(5.10), and
adding rational upper bounds gives \[
 \frac{92329284077596045529773584373060993024}
 {2507581295866705477237701416015625}<36821.
\] This proves (5.1). The moment calculation is finite and exact; it is
not an asymptotic assertion with an unspecified constant.

\subsubsection{5.4. Restoring the Gaussian model outside the
box}\label{restoring-the-gaussian-model-outside-the-box}

The exact leaf-coordinate moment-generating function under the
normalized real weighted Gaussian is \[
 \mathbb E e^{t\sqrt{n\beta}\eta_i}
 =e^{15t^2/32}\prod_{m=2,4,8,16}(1+t^2/m).                    \tag{5.11}
\] Indeed a leaf has four ancestor Haar coordinates with squared
coefficients \(1/2,1/4,1/8,1/16\). A union and Chernoff bound therefore
give \[
 P_{\mathcal B^c}\le32
 e^{-td\sqrt{19V/20}+15t^2/32}
 \prod_{m=2,4,8,16}(1+t^2/m).                                \tag{5.12}
\] The absolute value of (4.2) is at most \[
 Q/12+(V/18)S_4+(V^2/50)S_3^2.
\] Its weighted \(L^2\) norm is bounded by \(67/V\); the exact squared
constant is \[7405704307825/1693651716<67^2.\] Cauchy--Schwarz and the
conditioning factor in (3.5) bound the restored Gaussian-plus-correction
tail by \[
 2\{P_{\mathcal B^c}+(67/V)\sqrt{P_{\mathcal B^c}}\}|A_n(s)|.
\] For \(V\ge225\), choose \(t=5\) in (5.12). Exact rational exponential
and square-root bounds give \[
 P_{\mathcal B^c}<1/1000,\qquad
 2\{P_{\mathcal B^c}+(67/V)\sqrt{P_{\mathcal B^c}}\}<11/500.
                                                               \tag{5.13}
\] These are tails of the Gaussian model. They are not substituted for
the actual torus complement, which is bounded next.

\subsection{\texorpdfstring{6. The actual torus complement at
\(V\ge250\)}{6. The actual torus complement at V\textbackslash ge250}}\label{the-actual-torus-complement-at-vge250}

Represent each angle in \([-\pi,\pi]\); boundary sets have measure zero.
Call an angle central if \(|\alpha_i|\le\pi/3\). We divide the
complement of the two local boxes into five classes. All bounds in this
section are relative to the magnitude of \textbf{one} saddle,
\(|A_n(s)|\).

\subsubsection{6.1. Principal central sectors outside the
boxes}\label{principal-central-sectors-outside-the-boxes}

In the sector \(\sum\alpha_i=\pi\), set \(\eta_i=\alpha_i-\theta\), and
define \[
 g(x)=\frac{h_s(\theta)+h_s'(\theta)x-h_s(\theta+x)}{v_0}.
\] For either radial direction, the scalar inequalities imply at
\(d=2/5\) \[
 g(\pm d)\ge269/3750,\qquad
 \frac d{dr}g(\pm r)\big|_{r=d}\ge43/125.
\] These follow by integrating \(g''(0)\ge19/20\), \(|g'''(0)|\le1/4\),
and \(g''''\ge-3/2\). For larger radial distances still inside the
central angular interval, (3.4) gives \(g''\ge1/3\). Thus \[
 g(\pm r)\ge269/3750+(43/125)(r-d)+(r-d)^2/6\quad(r\ge d).
\] Let \[
 \gamma(R)=\frac13+\frac{158}{375R}-\frac{49}{625R^2}.
\] It is decreasing on \([2/5,5/4]\), and \(g(x)\ge\gamma(R)x^2/2\)
whenever \(|x|\le R\). The largest possible \(|\eta_i|\) in this sector
is \(19\pi/48<5/4\).

Split the maximum coordinate into the seventeen slabs \[
 [\ell_j,R_j]=[2/5+j/20,\,2/5+(j+1)/20],\qquad0\le j\le16.
\] On a slab, use \(K\le H\) and the Gaussian precision
\(\gamma(R_j)V\). Its normalization costs at most
\(2(17/(16\gamma(R_j)))^{23}\). Equation (5.11) gives the explicit slab
bound \[
 64\left(\frac{17}{16\gamma(R_j)}\right)^{23}
 e^{-t_j\ell_j\sqrt{\gamma(R_j)V}+15t_j^2/32}
 \prod_{m=2,4,8,16}(1+t_j^2/m).                               \tag{6.1}
\] At \(V=250\), take a rational lower bound for the square root by
flooring it to denominator \(1000\), and set \(t_j\) equal to the
integer part of the resulting \(\ell_j\sqrt{250\gamma(R_j)}\), minus
one, or one if this is smaller. Sum (6.1), bounding each negative
exponential by the reciprocal of a positive exponential Taylor partial
sum. The exact sum is less than \(9/80\). Each term decreases for larger
\(V\). Thus the two principal central complements together contribute
less than \(9/40\).

\subsubsection{6.2. Nonprincipal central
sectors}\label{nonprincipal-central-sectors}

The other central sectors have means \(\mu=\pm3\theta,\pm5\theta\).
Strong concavity gives \[
 \sum_i h_s(\alpha_i)
 \le16h_s(\mu)-\frac{v_0}{6}\sum_i(\alpha_i-\mu)^2.
\] The wide fourth-cumulant bound implies \[
 h_s(\theta)-h_s(3\theta)
 \ge v_0(4\theta^2-10\theta^4)>v_0/8.
\] For example, integrate \(\Re v(se^{ia})\ge v_0(1-3a^2/2)\) first from
zero to \(a\), then from \(\theta\) to \(3\theta\). The larger mean
\(5\theta\) only improves the loss. There are sixteen leaves, so the
total mean deficit exceeds \(2v_0\). The four sectors together are
bounded by \[
 8(51/16)^{23}e^{-2V}<10^{-10},\qquad V\ge250.                 \tag{6.2}
\]

\subsubsection{6.3. Exactly one noncentral
angle}\label{exactly-one-noncentral-angle}

Let its magnitude be \(a\in[\pi/3,\pi]\). From (3.1), \[
 h_s(a)\le-v_0\phi(a),\qquad
 \phi(a)=\frac{2-\cos a-\cos^2a}{3}.
\] For \(|x|\le2\pi/45\), (3.7) and Taylor's formula give \[
 h_s(x)\le-v_0x^2/2+v_0x^4/16\le-(49/100)v_0x^2.
\] The mean of the fifteen inliers has magnitude at least
\((\pi-a)/15\). Monotonicity of \(h_s\), followed by the short-angle
estimate at this smaller comparison argument, gives a mean cost at least
\((49/1500)v_0(\pi-a)^2\). The remaining inlier fluctuations have
Gaussian precision \(V/3\).

The function \(\phi(a)+(49/1500)(\pi-a)^2\) is bounded below by
\(5/12+49\pi^2/3375\). To check this, it is concave on
\([\pi/3,2\pi/3]\), because \(\phi''\le-1/6\) while the added second
derivative is \(49/750<1/6\); its first endpoint is the smaller one. On
\([2\pi/3,\pi]\), \(\phi\ge2/3\), which is larger than that endpoint
value. After subtracting (3.3), the deficit exceeds \[
 \delta_1=\frac{3246499}{13837500},                            \tag{6.3}
\] using \(\pi^2>493/50\).

Only the four ancestors of the outer leaf need their kernel factors
bounded by two. The eleven remaining contrasts lie orthogonally in the
fourteen-dimensional centered inlier space, and their product of
\(m_\nu/2\) is \(32\). Their Gaussian integral is
\(32(2\pi)^7(3/V)^{18}\). The original measure disintegrates as
\(da\,d\sigma/\sqrt{15}\). There are at most sixteen choices of leaf and
six winding sectors. After normalizing by (4.1), the total is bounded by
\[
 128\,3^{18}(17/16)^{23}V^{9/2}e^{-\delta_1V}<1/2000,
 \qquad V\ge250.                                             \tag{6.4}
\] The factor uses \(\sqrt{2\pi/15}<2/3\) and includes the conditioning
factor two.

\subsubsection{6.4. Exactly two noncentral
angles}\label{exactly-two-noncentral-angles}

The two leaf-to-root paths have at most seven distinct internal
ancestors. At least eight contrasts therefore involve only the fourteen
central leaves. Keep eight such contrasts and bound the other seven
kernel factors by two. Each retained node has size at most eight, so \[
 K\le2^{23}\prod_{j=1}^8x_j^2
\] in orthonormal coordinates of the thirteen-dimensional centered
inlier plane. Its Gaussian integral contributes
\((2\pi)^{13/2}(3/V)^{29/2}\). The coarea factor is \(1/\sqrt{14}\), and
there are at most \(\binom{16}{2}\cdot6=720\) leaf-pair/winding choices.

The deficit relative to the saddle is at least
\(2(5/12)-40/123=125/246\). Including the conditioning factor two, the
total bound is \[
 720\,2^{17}3^{14}(17/16)^{23}V^8e^{-125V/246}<10^{-10},
 \qquad V\ge250.                                             \tag{6.5}
\] Keeping the eight intact contrasts is essential: replacing the whole
kernel by a constant would introduce a much larger power of \(V\).

\subsubsection{6.5. At least three noncentral
angles}\label{at-least-three-noncentral-angles}

Here the deficit is at least \(3(5/12)-40/123=455/492\). Use
\(K\le2^{15}\) and normalized total torus volume one. The resulting
bound is \[
 128(44/7)^8(17/16)^{23}V^{23}e^{-455V/492}<10^{-10},
 \qquad V\ge250.                                             \tag{6.6}
\] No union multiplier is necessary: this is a pointwise bound on the
entire region.

All numerical inequalities (6.1)--(6.6) are checked by rational
arithmetic, positive exponential Taylor sums and elementary rational
bounds on \(\pi\). The functions \(V^a e^{-\delta V}\) in (6.4)--(6.6)
decrease for \(V\ge250\), as follows by differentiating their
logarithms.

\subsection{7. The uniform bulk theorem and phase
points}\label{the-uniform-bulk-theorem-and-phase-points}

Combining the two local errors, the two model-restoration errors, the
actual complement, and (4.4), we obtain \[
\begin{aligned}
\frac{|f_{n,4}(-s^{16})-2\Re A_n(s)|}{2|A_n(s)|}
 &<\frac{36821}{250^2}+\frac{53}{250}+\frac{11}{500}
 \\ &\quad+\frac12\left(\frac9{40}+\frac1{2000}+\frac3{10^{10}}\right)
 \\ &=\frac{18717720003}{20000000000}<\frac{47}{50}                 
\end{aligned}\tag{7.1}
\] whenever \(0<s\le1\) and \(V\ge250\). Notice that the actual
complement bounds were relative to one saddle; this accounts for the
factor \(1/2\) in the last term.

Define \[
 \Psi_n(s)=16n\arg p(\zeta)-\frac{45}{2}\arg v(\zeta).
\] Let \(m_q=zq'/q\). Direct algebra gives, for \(z=se^{i\theta}\), \[
 \Im m(z)=\frac{s\sin\theta(1+4s\cos\theta+s^2)}{|p(z)|^2}>0,
\] \[
 \Im m_q(z)=\frac{4s\sin\theta(1+s\cos\theta+s^2)}{|q(z)|^2}
             \le4\Im m(z).
\] The last inequality uses
\(|q|^2-|p|^2=6s\cos\theta(1+s^2)+15s^2\ge0\). Since \(zv'/v=1+m_q-2m\),
\[
 \frac{d\Psi_n}{d\log s}\ge(16n-45)\Im m(z)>0\qquad(n\ge3).  \tag{7.2}
\] Moreover, \[
 \Psi_n(0+)=-45\pi/32,\qquad \Psi_n(1)=n\pi.
\] Thus every integer phase level \(m\pi\), \(0\le m\le n\), has a
unique point in \((0,1]\). We use only levels \(m\ge250\), not formal
phase crossings near the endpoint.

For \(n\ge1001\), let \(s_*\) solve \(nv(s_*)=250\). The real variance
is strictly increasing on \((0,1)\), since \[
 sv'(s)=\frac{s(1-s^2)(1+7s+s^2)}{p(s)^3}>0.
\] The exact value \(v(7/20)=141260/346921\) shows \(s_*<7/20\). On this
interval, \[
 \frac{m(s)}{v(s)}-1=\frac{s(2s^2+2s-1)}{1+4s+s^2}\le0.
\] Also, for \(0\le\alpha\le\theta\), \[
 \Im v(se^{i\alpha})=
 \frac{s\sin\alpha(1-s^2)
 (4s^2\cos^2\alpha+8s^3\cos\alpha+8s\cos\alpha+s^4+5s^2+1)}
 {|p(se^{i\alpha})|^4}\ge0.
\] Therefore \(\arg p(se^{i\alpha})\) is concave in \(\alpha\), and
\(\arg p(\zeta)\le\theta m(s)\), while \(\arg v(\zeta)\ge0\). Hence \[
 \Psi_n(s_*)\le16n\theta v(s_*)=250\pi.                      \tag{7.3}
\] By monotonicity, all phase points \(s_{n,m}\), \(250\le m\le n\),
have \(V\ge250\). At such a point, \(\Re A_n=(-1)^m|A_n|\), so (7.1)
proves \[
 (-1)^m f_{n,4}(-s_{n,m}^{16})>0,
 \qquad250\le m\le n,\quad n\ge1001.                         \tag{7.4}
\] In particular, \(s_{n,n}=1\) and \(f_{n,4}(-1)\ne0\).

\subsection{8. The analytic edge family and exact enclosure
lemma}\label{the-analytic-edge-family-and-exact-enclosure-lemma}

\subsubsection{\texorpdfstring{8.1. A family regular at
\(u=0\)}{8.1. A family regular at u=0}}\label{a-family-regular-at-u0}

For complex \(u,z\) with \(|uz|<1\), define \[
 g_u(z)=\exp L(u,z),\qquad
 L(u,z)=\begin{cases}\log p(uz)/u,&u\ne0,\\z,&u=0.\end{cases} 
\] The logarithm is its analytic branch at zero. The displayed
singularity is removable. Set \[
 b_{0,k}(u)=[z^k]g_u(z),\qquad b_{j+1,k}=b_{j,k}^2-b_{j,k-1}b_{j,k+1},
 \qquad F_u(x)=\sum_{k\ge0}b_{4,k}(u)x^k.                     \tag{8.1}
\] These analytic functions are used only on domains where the bounds
below ensure convergence. For a reciprocal integer \(u=1/n\), \[
 b_{0,k}(1/n)=n^{-k}a_{n,k},\qquad
 F_{1/n}(x)=f_{n,4}(x/n^{16}).                               \tag{8.2}
\] The second identity follows because one application of \(\mathcal L\)
squares a geometric coefficient scaling. It is exact, not an asymptotic
limit.

\subsubsection{8.2. Complex-parameter
majorant}\label{complex-parameter-majorant}

The series \[
 \log p(w)=\sum_{\ell\ge1}c_\ell w^\ell,\qquad
 c_\ell=\begin{cases}1/\ell,&3\nmid\ell,\\-2/\ell,&3\mid\ell,\end{cases}
\] follows from \(p(w)=(1-w^3)/(1-w)\). Put \[
 G(q)=-\log(1-q)-\tfrac13\log(1-q^3),
 \qquad H(q)=\frac{qG'(q)-G(q)}{q^2}.
\] Let \(c\ge0\), \(r,R>0\), and \(q=(c+r)R<1\). On \(|u-c|\le r\),
\(|z|\le R\), \[
 |\partial_uL(u,z)|\le R^2H(q).
\] Integrating from \(c\) to \(u\) gives \[
 |g_u(z)|\le e^E,\qquad
 E=\frac{\log p(cR)}c+rR^2H(q),                              \tag{8.3}
\] where the quotient is interpreted as \(R\) when \(c=0\). At the real
starting parameter \(c>0\), the bound follows from \(|p(cz)|\le p(cR)\).

Cauchy's coefficient estimate followed by four applications of
\(\mathcal L\) yields \[
 |b_{4,k}(u)|\le2^{15}e^{16E}R^{-16k}.                       \tag{8.4}
\] Indeed the constants obey \(C_{j+1}=2C_j^2\), \(C_0=1\), while the
spatial exponent doubles at every layer. A missing negative-index
neighbor only improves the estimate.

\subsubsection{8.3. Moving-node Taylor
remainder}\label{moving-node-taylor-remainder}

Let a positive rational moving node on a real panel \([c-h,c+h]\) be \[
 t(u)=\frac{t_c}{1+a(u-c)},\qquad |a|r<1,\qquad0<h<r.
\] Then \(|t(u)|\le T=t_c/(1-|a|r)\) throughout the \textbf{complex}
parameter disk. Choose \(R>T\) with \((c+r)R<1\). For the truncated
evaluation \[
 Q_K(u)=\sum_{k=0}^K b_{4,k}(u)(-t(u)^{16})^k,
\] (8.4) gives \[
 |Q_K(u)|\le M=\frac{2^{15}e^{16E}}{1-(T/R)^{16}}.
\] If \(Q_K(c+hw)=\sum_{j\ge0}d_jw^j\), Cauchy's estimate proves \[
 \left|Q_K(c+hw)-\sum_{j=0}^{d}d_jw^j\right|
 \le M\frac{(h/r)^{d+1}}{1-h/r},\qquad |w|\le1.              \tag{8.5}
\] The finite coefficient truncation is bounded separately. For real
\(0\le u\le U\), \[
 \log p(uR_0)/u\le R_0+UR_0^2/2=:E_0,
\] because \(\log p(q)\le q+q^2/2\). If \(R_0>T_0=\max t(u)\) and
\(UR_0<1\), then \[
 |F_u(-t(u)^{16})-Q_K(u)|
 \le\frac{2^{15}e^{16E_0}(T_0/R_0)^{16(K+1)}}
             {1-(T_0/R_0)^{16}}.                            \tag{8.6}
\] The complex-disk radius in (8.5) need not equal the real
coefficient-tail radius in (8.6). In the certificates, \(r=4h\),
\(R=T+1\), and \(R_0=2T_0\) are used after checking the displayed domain
inequalities exactly.

\subsubsection{8.4. Exact finite Taylor arithmetic and the sign
test}\label{exact-finite-taylor-arithmetic-and-the-sign-test}

The retained Taylor coefficients are enclosed as \[
 d_i\in2^{e-P}[m_i-\rho,m_i+\rho],\qquad0\le i\le d,
\] with integer midpoints, a common nonnegative integer radius, and an
arbitrary integer binary exponent. Thus rounding contributes at most
\((d+1)\rho2^{e-P}\) on \(|w|\le1\).

For multiplication, exact integer convolution of the midpoint vectors
has raw coefficient error at most \[
 \rho_b\sum|m_a|+\rho_a\sum|m_b|+(d+1)\rho_a\rho_b.
\] Its raw exponent is \(e_a+e_b-P\). Rescaling a midpoint by a positive
integer denominator \(D\), with nearest integer \(n_i\), enlarges the
radius to \[
 \left\lceil\frac{\rho+\max_i|Dn_i-m_i|}{D}\right\rceil.
\] These formulas prove the enclosure invariant independently of
intermediate coefficient size. Normalization uses both midpoint
magnitudes and the radius; a zero midpoint with nonzero radius is not
discarded. For multiplication by a rational scalar \(N/D\), \(D>0\),
choose integer rounded midpoints \(n_i\); a valid new common radius is
\[
 \left\lceil\frac{|N|\rho+\max_i|Dn_i-Nm_i|}{D}\right\rceil.
\] Addition first aligns binary exponents by these outward-rounded
rescalings and then adds the radii. A formal coefficient shift preserves
the common radius on retained entries. Thus all primitive operations
used in the recurrence preserve the stated coefficientwise enclosure.

The initial jets are generated from the exact recurrence \[
 b_{0,0}=b_{0,1}=1,\qquad
 (k+1)b_{0,k+1}=(1-ku)b_{0,k}+u(2-(k-1)u)b_{0,k-1}.           \tag{8.7}
\] Four radius-one stencils give the shared \(b_{4,k}\) jets. At each
moving node, Horner evaluation uses the Taylor jet \[
 -t_c^{16}(1+ahw)^{-16},\qquad
 [w^j]=-t_c^{16}\binom{15+j}{j}(-ah)^j.
\] Truncating after formal degree \(d\) cannot affect retained formal
coefficients.

Convert the midpoint polynomial to Bernstein form after \(w=2x-1\),
\(0\le x\le1\). Its minimum and maximum controls enclose all values. Add
the rounding envelope and the two independent errors (8.5)--(8.6). A
strictly positive lower bound or strictly negative upper bound then
certifies the sign on the entire closed panel. This is an interval-wide
proof, not a sampled evaluation.

All transcendental majorants can be replaced by rational upper bounds.
For example, expand \(H(q)\) through degree thirty and bound the
remainder by \(2q^{31}/(1-q)\); bound logarithms by the positive atanh
series with its geometric tail; use a positive rational lower bound for
\(\log2\) to replace \(e^X\) by
\(2^{\lceil X/\log2_{\rm lower}\rceil}\). Consequently the final sign
comparison is an exact rational inequality.

\subsection{9. Parameter-wide bulk bypass and the final
join}\label{parameter-wide-bulk-bypass-and-the-final-join}

For real \(r>0\), the analytic function \(p(z)^r\) on \(|z|<1\) is
defined with the logarithm at zero. The diagonal identity, kernel
symmetrization and all modulus estimates remain valid: integrality of
\(r\) was not used in those steps. Thus (7.1) extends to real \(r\),
\(0<s<1\), \(rv(s)\ge250\). No polynomial degree or reciprocal-root
assertion is made for noninteger \(r\).

At a moving node, set \(r=1/u\), \(s=ut(u)\). If \[
 rv(s)\ge250,\qquad |\Psi_r(s)-j\pi|\le1/3,
\] then (7.1) gives the desired sign because
\(\cos(1/3)\ge17/18>47/50\). At \(u=0\), the analytic family tends
uniformly on compact sets to its Poisson-scale limit \(g_0(z)=e^z\). The
amplitude has a finite positive limiting magnitude; the strict uniform
margin therefore extends the sign to that endpoint.

These bypass inequalities and the final join are themselves low-degree
rational polynomial certificates. Write \(t=A/B\), \(B>0\), and let
\(c=\cos\theta,d=\sin\theta,c_2=\cos2\theta,d_2=\sin2\theta\). Then \[
 \arg p(ut e^{i\theta})=\arctan\frac{uN_p}{D_p},
\quad D_p=B^2+cuAB+c_2u^2A^2,
\quad N_p=dAB+d_2uA^2,
\] with \(D_q=B^2+4cuAB+c_2u^2A^2\), \(N_q=4dAB+d_2uA^2\) for \(q\).
Hence \[
 \Psi_{1/u}(ut)
 =(16/u+45)\arctan(uN_p/D_p)
 -(45/2)\arctan(uN_q/D_q)-45\pi/32.                          \tag{9.1}
\] All real parts are positive, so the branches are unambiguous. Use
certified rational sine, cosine and \(\pi\) enclosures, together with \[
 x-x^3/3\le\arctan x\le x-x^3/3+x^5/5\qquad(x\ge0).
\] The upper polynomial is increasing; the lower is increasing on
\([0,1]\), which is verified for the relevant ratio bounds. After
division by \(u\), the apparent singularity in (9.1) is removable.
Clearing positive denominators gives degree-sixteen polynomial phase
margins. The variance condition is the degree-four inequality obtained
from \[
 \frac{v(ut)}u=
 \frac{AB(B^2+4uAB+u^2A^2)}{(B^2+uAB+u^2A^2)^2}.              \tag{9.2}
\] Positive Bernstein controls certify both types of margin on each
whole panel.

For the supplied final harmonic node, the exact inequalities establish
\(\Psi_{1/u}(ut_{249}(u))<250\pi\). For the root argument we use this
only at reciprocal integers \(u=1/n\): phase monotonicity then gives \[
 t_{249}(1/n)<n s_{n,250}.                                   \tag{9.3}
\]

\subsection{10. Certified edge inputs and
assembly}\label{certified-edge-inputs-and-assembly}

The rational parameter grid is \[
 u_\ell=\ell/32032,\qquad0\le\ell\le32.
\] For each \(j=1,\ldots,249\), the endpoint table supplies positive
rational values \(T_{j,\ell}\). On a panel interpolate reciprocals: \[
 t_j(u)=\left(\frac{1-\lambda}{T_{j,\ell}}
               +\frac{\lambda}{T_{j,\ell+1}}\right)^{-1},
 \quad \lambda=\frac{u-u_\ell}{u_{\ell+1}-u_\ell}.
\] The exact adjacent endpoint comparisons imply strict node ordering
throughout every panel: reciprocal interpolation preserves the reversed
inequalities, and the final reciprocal restores their order.

The completed finite edge certificate consists of panels \(0,\ldots,7\),
exactly covering \[
 [0,8/32032]=[0,1/4004].
\] On each of these panels, every node \(j=1,\ldots,249\) has the strict
sign \[
 (-1)^jF_u(-t_j(u)^{16})>0.                                  \tag{10.1}
\] The certificate inputs are as follows; the entries distinguish Taylor
enclosures from the exact bulk-bypass criterion, not separate
mathematical conclusions.

\begin{longtable}[]{@{}rrrr@{}}
\toprule\noalign{}
Panel & Taylor-enclosed signs & Bulk-bypass signs & Total \\
\midrule\noalign{}
\endhead
\bottomrule\noalign{}
\endlastfoot
0 & 246 & 3 & 249 \\
1 & 245 & 4 & 249 \\
2 & 243 & 6 & 249 \\
3 & 241 & 8 & 249 \\
4 & 240 & 9 & 249 \\
5 & 239 & 10 & 249 \\
6 & 238 & 11 & 249 \\
7 & 237 & 12 & 249 \\
\end{longtable}

Each Taylor entry includes an exact final coefficient enclosure and
strict rational comparison against both analytic errors. The hypotheses
of (8.3)--(8.6) are checked entry by entry. The bypass entries use the
exact inequalities in §9. The final join (9.3) and ordering use the same
frozen rational node table. None of these statements follows merely from
the aggregate counts in the table.

Now fix an integer \(n\ge4004\). Its reciprocal lies in the certified
closed parameter union. Since \(F_{1/n}(0)=1\), the first 249
alternating edge signs give 249 disjoint root intervals. The last edge
sign is negative, while phase point 250 has positive sign, so (9.3)
supplies one more interval. Bulk phases 250 through \(n\) supply
\(n-250\) further intervals. In the original polynomial variable these
are \[
 249+1+(n-250)=n
\] disjoint intervals inside \((-1,0)\). Thus \(f_{n,4}\) has at least
\(n\) distinct roots there.

The operator \(\mathcal L\) preserves palindromicity, degree \(2n\), and
both endpoint coefficients one. Reciprocal roots give another \(n\)
distinct roots below \(-1\). The value at \(-1\) is nonzero by (7.4).
These \(2n\) roots exhaust the degree, so every root is strictly
negative and real.

The companion finite/algebraic chapter proves nonnegativity of iterates
zero through three for \(n\ge56\), by finite PF4 seed certificates and
convolution closure. This includes the entire infinite range used here.
Brändén's theorem then gives all iterates from four onward. For
\(10\le n\le4003\), exact directed ratio certificates enter the
invariant cone \(\max q_k\le1/3\), which implies all subsequent iterates
are positive. The finite and infinite ranges meet without a gap and
cover every integer \(n\ge10\).

\subsection{11. Reproducibility
boundary}\label{reproducibility-boundary}

The mathematical proof requires more than a floating plot or a list of
successful samples. Its finite data are:

\begin{itemize}
\tightlist
\item
  the rational Bernstein certificates for the scalar inequalities in §3;
\item
  the exact weighted-moment sum in §5;
\item
  the rational complement bounds in §6;
\item
  the frozen moving-node endpoint table, interval-wide sign enclosures,
  and final-join controls;
\item
  the finite cone and PF4 certificates described in the companion
  chapter.
\end{itemize}

The computation uses reviewed exact integer/rational arithmetic and
outward-rounded coefficient enclosures. An executable certificate
checker must verify domains, strict sign inequalities, index coverage
and the connection between every saved enclosure and its parameters.
Cryptographic hashes help identify the intended records; hashing alone
does not establish any mathematical inequality. This draft presents a
computer-assisted proof, not a claim of formal proof-assistant
verification or external peer review.

\subsubsection{Primary references}\label{primary-references}

\begin{enumerate}
\def\labelenumi{\arabic{enumi}.}
\tightlist
\item
  L. A. Medina and A. Straub, \emph{On multiple and infinite
  log-concavity}, 2014 preprint; Annals of Combinatorics \textbf{20}
  (2016), 125--138.
  \href{https://arminstraub.com/downloads/pub/infinitelogconcavity.pdf}{Author-hosted
  paper}, especially Example 5.3 and Table 1.
\item
  P. Brändén, \emph{Iterated sequences and the geometry of zeros},
  Journal für die reine und angewandte Mathematik \textbf{658} (2011),
  115--131. \href{https://arxiv.org/pdf/0909.1927}{arXiv version},
  especially §3.
\end{enumerate}

\clearpage

\section{Reproducing the computational
premises}\label{reproducing-the-computational-premises}

The accompanying \texttt{certificate\_package} has three independent
directories: \texttt{finite}, \texttt{analytic}, and \texttt{edge}. It
requires Python with assertions enabled. Exact analytic polynomial
checks use SymPy 1.14.0, and accelerated edge regeneration uses
python-flint 0.9.0. The finite checks and saved-edge replay require only
the standard library. Dependency versions and backend provenance are
pinned in the package.

From the package directory, the following commands reconstruct the
stated premises:

\begin{Shaded}
\begin{Highlighting}[]
\ExtensionTok{python}\NormalTok{ finite/verify\_pf4.py}
\ExtensionTok{python}\NormalTok{ finite/verify\_finite.py }\AttributeTok{{-}{-}full} \AttributeTok{{-}{-}output}\NormalTok{ finite\_result.json}
\ExtensionTok{python}\NormalTok{ analytic/run\_analytic.py }\AttributeTok{{-}{-}full}
\ExtensionTok{python}\NormalTok{ edge/regenerate.py }\AttributeTok{{-}{-}full}
\ExtensionTok{python}\NormalTok{ edge/replay\_saved\_inequalities.py}
\end{Highlighting}
\end{Shaded}

The finite command verifies all 3994 exponents in the finite range, not
a selection. The analytic command reconstructs the rational polynomial
and weighted-moment arithmetic underlying the displayed inequalities.
The edge-generation command reconstructs all interval enclosures from
the coefficient recurrence and four applications of the operator.
Finally, the saved-inequality command evaluates every enclosure's strict
sign margin and its two analytic remainder budgets.

For a shorter inspection,
\texttt{python\ edge/replay\_saved\_inequalities.py} alone checks all
1929 saved enclosure inequalities, 63 bypass boxes, node ordering and
the final join. This shorter command is deliberately described as
\emph{conditional on the saved enclosures}: it does not reproduce their
interval-propagation provenance. The separate representative command

\begin{Shaded}
\begin{Highlighting}[]
\ExtensionTok{python}\NormalTok{ edge/regenerate.py }\AttributeTok{{-}{-}panel}\NormalTok{ 0 }\AttributeTok{{-}{-}node}\NormalTok{ 200}
\end{Highlighting}
\end{Shaded}

reconstructs one enclosure and compares every deterministic midpoint,
radius and exponent. A successful representative test is not represented
as a full-chain rerun. The original complete generation and independent
saved-inequality checks were performed; the portable full-generation
command is supplied for reproduction.

Eight deterministic compressed data files contain the original final
jets, one file per parameter panel. Hash manifests authenticate the
files and identify the exact reproducibility version. Hash agreement is
not a substitute for the mathematical implications proved in the
manuscript or for arithmetic reconstruction. Intermediate shared
coefficient checkpoints are omitted from the compact distribution; full
regeneration rebuilds them in memory.

The computational evidence and proof composition have received
independent internal checks. External peer review and proof-assistant
formalization remain distinct activities; neither is claimed by the
present draft.

\clearpage

\section{References}\label{references-1}

\begin{enumerate}
\def\labelenumi{\arabic{enumi}.}
\tightlist
\item
  L. A. Medina and A. Straub, \emph{On multiple and infinite
  log-concavity}, arXiv:1405.1765v1 (2014); \emph{Annals of
  Combinatorics} \textbf{20} (2016), 125--138. The specific target is
  Example 5.3, printed page 11 of the preprint.
  \url{https://arxiv.org/abs/1405.1765v1}
\item
  P. Brändén, \emph{Iterated sequences and the geometry of zeros},
  arXiv:0909.1927v3 (2010); \emph{Journal für die reine und angewandte
  Mathematik} \textbf{658} (2011), 115--131. Section 3 gives the
  real-nonpositive-zero preservation used after the fourth iterate.
  \url{https://arxiv.org/abs/0909.1927v3}
\item
  P. R. W. McNamara and B. E. Sagan, \emph{Infinite log-concavity:
  developments and conjectures}, \emph{Advances in Applied Mathematics}
  \textbf{44} (2010), 1--15; arXiv:0808.1065. The symmetric-function
  iteration identities and Toeplitz positivity provide context for the
  earlier-layer argument. \url{https://arxiv.org/abs/0808.1065}
\item
  I. G. Macdonald, \emph{Symmetric Functions and Hall Polynomials},
  second edition, Oxford University Press, 1995, Chapter I, Sections 3,
  5 and 9. Used for dual Jacobi--Trudi and Littlewood--Richardson
  identities.
\end{enumerate}

The exact polynomial-multiplication implementation uses the released
python-flint 0.9.0 interface to FLINT. This is an arithmetic dependency,
not a mathematical theorem proved here. Its release provenance and exact
conversion contract are documented with the verification software.
\url{https://pypi.org/project/python-flint/0.9.0/}

\end{document}
